\origpage{325}{N. Robertson, P.D. Seymour / Journal of Combinatorial Theory, Series B 92 (2004) 325--357}
{\centering
{\Large Graph Minors. XX. Wagner's conjecture\par}\medskip

Neil Robertson$^{\mathrm{a}, 1}$, P.D. Seymour$^{\mathrm{b}, 2}$\par
$^{\mathrm{a}}$\textit{Department of Mathematics, Ohio State University, 231 West 18th Ave., Columbus, OH 43210, USA}\par
$^{\mathrm{b}}$\textit{Telcordia Technologies, 445 South St., Morristown, NJ 07960, USA}\par
Received 13 February 2001\par
}
\begin{quote}
\textbf{Abstract}

We prove Wagner's conjecture, that for every infinite set of finite graphs, one of its members is isomorphic to a minor of another.

\textit{Keywords:} Graph; Minor; Surface embedding; Well-quasi-ordering
\end{quote}

\sectionhead{1. Introduction}

A famous conjecture of Wagner [6] asserts that for any infinite set of graphs, one of its members is isomorphic to a minor of another (all graphs in this paper are finite). It has been one of the main goals of this series of papers to prove the conjecture, and in this paper the proof is completed.

Our method is roughly as follows. If $\{G_1, G_2, \ldots\}$ is a counterexample to Wagner's conjecture then none of $G_2, G_3, \ldots$ has a minor isomorphic to $G_1$, and so to prove Wagner's conjecture it suffices to show the following.

\textbf{1.1.} \textit{For every graph $H$ and every infinite set of graphs each with no minor isomorphic to $H$, some member of the set is isomorphic to a minor of another member of the set.}%
\origfootnote[2]{This research was partially supported by NSF Grant DMS 8504054 and was partially performed under a consulting agreement with Bellcore.}%
\origfootnote[3]{Present address: Mathematics Department, Princeton University, 201 Fine Hall, Princeton, NJ 08544, USA.}
